\section{Data and Evaluation Protocol}
\label{sec:data}

\subsection{CATNAP cohort and labels}

We compile a dataset from the CATNAP database, which aggregates antibody
sequences and quantitative viral neutralizations. The half-maximal
inhibitory concentration ($\text{IC}_{50}$) measures an antibody's potency
against a specific viral isolate: it gives the concentration of antibody
needed to neutralize 50\% of the virus (lower values indicate
higher antibody potency). 
We evaluate each antibody--virus as neutralized,
if it is neutralized at the resistance threshold of $50\,\mu\mathrm{g/mL}$. 
Neutralization breadth
$b_i$, for an antibody $i$, is the ratio of of tested antibody--virus
measurements neutralized below the threshold, also given here:
\begin{equation}
    b_i = \frac{\#\{\text{measurements with IC}_{50}\leq
    50\,\mu\mathrm{g/mL}\}}{\#\{\text{tested antibody--virus measurements}\}},
\end{equation}
We restrict our analysis to antibodies tested against at least 20 virus
strains for decreased variance in selection. For a binary classification
task, we assign the "broad" label ($1$) where $b_i\geq0.50$, amd the narrow/weak label
($0$) where $b_i\leq0.20$. We exclude the middle band to establish
a distinct boundary between the binary classes.
[cite a reason for the choice of the middle band region, size]
We check structural variable-heavy-region ($V_H$)
completeness by requiring conserved cysteine residues and framework
markers to isolate the variable domain[explain this, and CDR in the intro.]
(Figure~\ref{fig:antibody_structure} shows the $V_H$ domain
within the overall antibody structure)
while retaining somatic mutations.

\begin{figure}[t]
  \centering
  \includegraphics[width=0.8\linewidth]{figures/antibody_structure.png}
  \caption{Schematic diagram of IgG antibody structure, highlighting the heavy and light chains, variable and constant domains, antigen-binding site, and disulfide bonds (reproduced from Mamahdi14~\cite{Mamahdi2016Antibody}).}
  \label{fig:antibody_structure}
\end{figure}

The final labeled set contains the 300 unique $V_H$ sequences (202 broad and 98
weak), with sequence lengths ranging from 116 to 146 amino acids.
For continuous-breadth regression, we predict the ratio $b_i$ directly for
the 300 selected assays.

\subsection{Lineage grouping and split}

To prevent data leakage of closely related sequences,
we group antibodies by distinct lineages, we do a robust 
hierarchical sorting to
ensure lineage stratificaton. The hierarchical sort uses
CATNAP clonal-lineage metadata (216 antibodies), patient
metadata indicating the same donor (70), sequence-level clustering via
single-linkage clustering at a threshold of $\geq 80\%$ Levenshtein sequence
identity (2 antibodies), and the remaining individual singletons (12). This
produces 113 unique group IDs. The tiers are not cross-merged, leaving nine
near-duplicate pairs that span a group boundary, which we audit but
intentionally do not merge (Appendix~\ref{app:grouping_audit}).
Across 774 in-group pairs, median pairwise
sequence identity is $0.779$ (Interquartile Range [IQR] $0.704$--$0.885$);
across 44,076 out-group pairs, the median is $0.424$ (IQR $0.377$--$0.504$).

Both classification and regression tasks utilize a single 5-fold
\texttt{StratifiedGroupKFold} split. This partitioning scheme guarantees that
all antibodies belonging to any of the 113 lineage groups remain together in
either the training or the validation set (unseen group evaluation), while
simultaneously stratifying folds to maintain the overall binary label ratio
(broad vs. narrow) in each split.

\subsection{Metrics, selection, and uncertainty}

We evaluate classification performance using the Area Under the Receiver
Operating Characteristic curve (AUROC) and Area Under the Precision-Recall Curve
(AUPRC, or average precision). AUROC measures the probability that the
classifier ranks a random positive instance higher than a random negative one.
AUPRC evaluates the precision-recall trade-off, this is informative
under class imbalances. For regression on the continuous breadth metric,
we report: the
Pearson correlation coefficient ($r$), which measures linear correlation; the
Spearman rank correlation ($\rho$), which evaluates monotonic but non-linear
relationships; and the coefficient of determination ($R^2$), representing the
proportion of target variance explained by the model compared to a simple mean
baseline.

For hyperparameter tuning and configuration selection, we use the mean
score
across the 5 validation folds (mean AUROC for classification; mean Spearman
$\rho$ for regression). After selecting a configuration, we report the primary
performance metrics on pooled out-of-fold (OOF) predictions
(concatenating the held-out predictions from all folds). This pooling prevents
bias from differing fold sizes and class compositions. To estimate prediction
variance, we calculate 95\% uncertainty intervals using 10,000 bootstrap draws,
resampling complete lineage groups with replacement rather than individual
sequences. To compute paired deltas that evaluate model differences, we
resample the groups and calculate the metric difference on aligned OOF
predictions. These intervals quantify sensitivity to lineage composition but do
not account for configuration-selection bias.

\subsection{OAS adaptation corpus}

The executed domain-adaptation run uses 219,837 globally deduplicated human
heavy-chain sequences from 11 OAS data units (108,675 clonotypes), with the
largest unit contributing 20,000 sequences (9.1\%). Whole data units that are
systematically 5$'$ truncated do not enter the corpus. A clonotype-aware split produces
198,130 training and 21,707 validation sequences. OAS supplies repertoire data
for adaptation only; we neither label OAS sequences as HIV bnAbs nor include them
in downstream test folds~\cite{olsen2022oas}.

\subsection{Lineage-interpolation cohort}

The interpolation track unions the 300 labeled antibodies with 11 trajectory
rows; four trajectory IDs already occur in the labeled cohort, yielding 307
unique sequences and 39,129 residue embeddings. The VAE training split contains
214 antibodies (27,490 residues), a group-held validation split contains 46
antibodies (5,738 residues), and the complete VRC34, VRC01, DH270, and DH511
lineages form the holdout set (47 antibodies; 5,901 residues).
